\clearpage
\chapter*{강연 212의 정오표}
\addcontentsline{toc}{chapter}{강연 212의 정오표}
\setcounter{footnote}{2}
\src{301}

\noindent\textbf{5. TDTE III: 몫 준스킴.}

\begin{center}
[부르바키 세미나, 제13권, 1960/61, 제212호, 20쪽.]
\end{center}

\noindent\textbf{15쪽, 주장 (i).} --- $Y=X/R$가 $S$ 위에서
준사영이라는 사실은, 현재로서는 $R$이 동치관계에서 유래하는 경우에만
증명된 것으로 보인다.

\noindent\textbf{18쪽과 19쪽.} --- 다음 강연의 말미에서 지적하듯이,
문제의 추측은 분명히 거짓이다. 주 8.1에서 언급한 ``긍정적인 사실''은
여러 저자(나가타, 로젠리히트, 그로텐디크, \ldots)가 동시에 증명한 것으로
보인다.

TDTE V에서 피카르 스킴을 구성하기 위해 여기서 전개한 몫을 취하는 이론을
적용한 부분은, 힐베르트 스킴을 적절히 이용하는 것으로도 대체할 수 있음을
지적해 둔다(cf. SHMT). § 8에서 말했듯이, 여기서 제시된
\src{302}
이론의 가장 큰 공백은 고유하지 않은 동치관계로 나눈 몫의 존재 판정법이
없다는 점이다. 그러한 동치관계에는 사영군의 어떤 작용들에서 유래하는
것들이 있다. 이 방향에서 중요한 정리 하나를 멈퍼드가 얻었다\footnote{MUMFORD
(David). --- An elementary theorem in geometric invariant theory, \emph{Bull.
Amer. Math. Soc.}, t. 67, 1961, p. 483--486.}. 그의 결과를 더 정교하게 만든
형태와 이 이론에 대한 여러 응용에 대해서는 SHMT를 보라.
